% GENERATED FILE. Edit canonical lessons, then run npm run generate:books.
% canonical-source-hash: fnv1a64:b304322e206ddaf6
% canonical-lessons: PT-C263-desculpe, PT-C263-direita-esquerda, PT-C263-revisao-caminho

\chapter{Desculpe, Como Se Vai Para --- Asking the Way}
\label{ch:pt-asking-the-way}
\hlchaptermodality{263}

\begin{chapteropening}
\textbf{By the end of this chapter:} \emph{I can stop a stranger politely, ask the way to a place, and follow an answer with à direita, à esquerda and em frente.}

Complete PT-C263-revisao-caminho: ask the way, follow the answer, check it back and thank.
\end{chapteropening}

\section[desculpe, como se vai para…?]{desculpe, como se vai para… — the question that finds the way}

\label{lesson:PT-C263-desculpe}



\begin{quote}
Say \textquotedblleft{}to go\textquotedblright{}. (\emph{Ir.}) Say \textquotedblleft{}the station\textquotedblright{}. (\emph{A estação.})

You are in a town you do not know, and the station is somewhere.
\end{quote}

\subsection*{You'll want to know: desculpe}
\textbf{desculpe} — \textquotedblleft{}excuse me\textquotedblright{}, said to someone you do not know.

It is the polite form of \textbf{desculpar}, \textquotedblleft{}to excuse\textquotedblright{}, the one that goes with \emph{o senhor}, \emph{a senhora} and \emph{você}, which a stranger in the street is owed. The same word apologises: bump into someone and you say \emph{desculpe!} In Brazil you will also hear \emph{com licença}, \textquotedblleft{}with permission\textquotedblright{}, for this; in Portugal you hear that mostly when someone needs to get past.

\begin{grammarlens}[title={Grammar Lens: como se vai para…?}]
\begin{quote}
\textbf{Desculpe, como se vai para a estação?} — Excuse me, how do I get to the station?
\end{quote}

Word for word, \textquotedblleft{}how does one go to the station\textquotedblright{}. The \textbf{se} leaves the person out, so the question is about the way, not about you. \textbf{para} plus the place, with its article:

\begin{quote}
\textbf{para a} estação · \textbf{para o} museu · \textbf{para o} hotel
\end{quote}

A second way to ask uses \textbf{fica}, the verb that tells where a place is: \emph{onde fica a estação}, where the station is. \emph{Como se vai para} asks for the \textbf{way}; \emph{onde fica} asks for the \textbf{place}. Either can be followed by a short question that saves a walk, \emph{é longe}, is it far.
\end{grammarlens}

\begin{cousinweb}
\textbf{desculpar} is built on Latin \emph{culpa}, \textquotedblleft{}fault, blame\textquotedblright{} (Lewis \& Short), with \emph{des-}, \textquotedblleft{}away\textquotedblright{}, in front: to excuse someone is to take the fault away from them. EtymDB 2.1 has no row for \emph{desculpar} itself, so that is the usual reading of its parts. English \textbf{culpable} shares the \emph{culpa}.
\end{cousinweb}

\subsection*{Your turn}
Say these aloud:

\begin{itemize}
\raggedright
  \item \emph{Desculpe, como se vai para a estação?} — and let the voice rise
  \item the same, to the museum, then to the hotel
  \item ask for the place instead, and how far — \emph{Desculpe, onde fica a estação? É longe?}
\end{itemize}

\begin{tcolorbox}[breakable,colback=teal!4,colframe=teal!35!black,title={Before you move on}]
Who gets \emph{desculpe}? (A \textbf{stranger}, anyone you call \emph{o senhor}, \emph{a senhora} or \emph{você}.) Ask the way to the hotel. (\emph{\textbf{Desculpe, como se vai para o hotel?}}) Which verb asks for the place rather than the way? (\emph{\textbf{Fica}} — \emph{onde fica}.)
\end{tcolorbox}
\section[à direita, à esquerda]{à direita, à esquerda — following the answer}

\label{lesson:PT-C263-direita-esquerda}



\begin{quote}
Stop a stranger and ask the way to the station. (\emph{Desculpe, como se vai para a estação?})

Asking was the easy half. Now the answer comes back, fast.
\end{quote}

\subsection*{You'll want to know: à direita, à esquerda}
\textbf{à direita} — \textquotedblleft{}on the right, to the right\textquotedblright{}. \textbf{à esquerda} — \textquotedblleft{}on the left, to the left\textquotedblright{}. With \textbf{em frente}, \textquotedblleft{}straight ahead\textquotedblright{}, which you already have, that is every direction a street can send you.

\begin{quote}
\emph{A estação é à direita.} — The station is on the right. \emph{A primeira à esquerda.} — The first on the left.
\end{quote}

That second one is a whole answer in four words: \emph{a primeira} agrees with \emph{a rua}, the street it leaves unsaid. The accent on \textbf{à} marks two words in one, \emph{a}, \textquotedblleft{}to\textquotedblright{}, and \emph{a}, \textquotedblleft{}the\textquotedblright{}, run together.

\begin{grammarlens}[title={Grammar Lens: siga em frente, vire à direita}]
A stranger answering a \emph{desculpe} answers politely, and the directions are instructions:

\begin{quote}
\textbf{Siga} em frente e \textbf{depois vire} à direita no semáforo. — Go straight on, and then turn right at the traffic lights.
\end{quote}

\textbf{vire} is \textbf{virar}, \textquotedblleft{}to turn\textquotedblright{}, made a polite instruction. \textbf{siga} is the same for \textbf{seguir}, \textquotedblleft{}to go on\textquotedblright{}, the verb inside your \emph{a seguir}. You need to \textbf{understand} these two, not build them yet.

A long answer comes in the order you will walk it, \emph{depois} marking each piece.
\end{grammarlens}

\begin{cousinweb}
\textbf{direita} has no parent in EtymDB 2.1. It is usually traced to Latin \emph{dīrēctus}, \textquotedblleft{}straight, set in a straight line\textquotedblright{}, the participle of \emph{dīrigere} (Lewis \& Short, under \emph{dirigo}).

\textbf{esquerda} is the surprise. It comes from Basque \emph{ezker}, \textquotedblleft{}left\textquotedblright{} (EtymDB 2.1), a language spoken in Iberia before Latin arrived, and so does Spanish \emph{izquierda}. Latin's own word for the left hand, \emph{sinistra}, lives on in English \textbf{sinister}.
\end{cousinweb}

\subsection*{Your turn}
Say these aloud:

\begin{itemize}
\raggedright
  \item \emph{à direita}, \emph{à esquerda}, \emph{em frente}
  \item \emph{A primeira à esquerda.} — then \emph{A segunda à direita.}
  \item the stranger's answer — \emph{Siga em frente e depois vire à direita no semáforo.}
  \item say the answer back to check it — \emph{Em frente, depois à direita no semáforo?}
\end{itemize}

\begin{tcolorbox}[breakable,colback=teal!4,colframe=teal!35!black,title={Before you move on}]
Which word is \textquotedblleft{}left\textquotedblright{}, and where does it come from? (\emph{\textbf{Esquerda}} — from \textbf{Basque}.) What is \emph{vire}, and from which verb? (A polite \textquotedblleft{}turn\textquotedblright{}, from \emph{\textbf{virar}}.) What does \emph{a primeira à direita} leave unsaid? (\textbf{The street} — \emph{a primeira} agrees with it.)
\end{tcolorbox}
\section[(review)]{Review — ask, listen, check, thank}

\label{lesson:PT-C263-revisao-caminho}



\begin{quote}
Say \textquotedblleft{}on the left\textquotedblright{}, then \textquotedblleft{}on the right\textquotedblright{}. (\emph{À esquerda, à direita.})
\end{quote}

\begin{grammarlens}[title={What you've built}]
Asking the way has four steps, and you now have Portuguese for every one.

\textbf{Ask}: \emph{Desculpe, como se vai para a estação?} \textbf{Listen}: \emph{Siga em frente e depois vire à esquerda no cruzamento.} \textbf{Check}, by saying the answer back as a question: \emph{Em frente, depois à esquerda?} \textbf{Thank}: \emph{Obrigado!}, or \emph{Obrigada!}

The check is the step travellers skip, and it is the one that saves the walk. A stranger who hears their own directions said back will correct the one piece you got wrong. And every \emph{depois} in the answer is Latin \emph{post}, \textquotedblleft{}after\textquotedblright{}: one step, then the next.
\end{grammarlens}

\subsection*{Your turn}
Say these aloud:

\begin{itemize}
\raggedright
  \item ask — \emph{Desculpe, como se vai para o museu?}
  \item hear the answer and say it back — \emph{Siga em frente e vire à direita.} … \emph{Em frente, depois à direita?}
  \item thank — \emph{Obrigado!} or \emph{Obrigada!}
  \item now be the stranger: send someone to the station — straight ahead, then the first on the left
\end{itemize}

\begin{tcolorbox}[breakable,colback=teal!4,colframe=teal!35!black,title={Before you move on}]
What are the four steps? (\textbf{Ask}, \textbf{listen}, \textbf{check}, \textbf{thank}.) Which one do travellers skip? (\textbf{The check.}) If a friend is lost, advise them to ask, and explain why. (\emph{Acho que \textbf{devias} perguntar, porque a estação é longe.}) And what does \emph{explicar} do, word for word? (It \textbf{unfolds}.)
\end{tcolorbox}
